\section{Motivation and requirements}\label{sec:motivation}

The statistical properties of free-flight trajectories (time spent in thermals, length of the
glides, mean squared displacement) are usually estimated by treating flights as independent
samples. The assumption fails when several pilots fly together: they share thermals and routes,
and their trajectories are correlated. Group flight may also change the behaviour itself. In
soaring birds, social information on the position of thermals changes the gliding
strategy~\cite{williams2018}, and in a migrating flock a few individuals find the thermals while
the others follow them~\cite{flack2018,nagy2018}. Comparing group flights with solo flights
therefore needs a rule that says, for every flight and every instant, whether the pilot is in a
group and with whom.

The intuitive notion of \emph{flying together} that the rule must make operational is the
following (Figure~\ref{fig:requirements}).
\begin{enumerate}[label=(R\arabic*), leftmargin=3.2em]
  \item The pilots take off from the same area, not necessarily from the same takeoff.
  \item Departures may be staggered, but not to the point of losing sight of each other.
  \item The pilots follow roughly the same route: nearby places and similar motion, not
        necessarily the same points.
  \item The composition of the group changes in time: some pilots join, others change route.
\end{enumerate}

\begin{figure}
\centering
\begin{tikzpicture}[x=0.97cm,y=1cm,
    tag/.style={draw=ink2, rounded corners=2pt, fill=white, font=\footnotesize\bfseries, inner sep=2pt}]
  \draw[light, line width=18pt, line cap=round] plot[smooth] coordinates {(2.6,2.35) (6,2.75) (9,2.4) (12,2.4) (14.6,2.4)};
  \draw[ink2, densely dashed] (1.15,2.25) ellipse (0.95 and 1.25);
  \fill[ink] (0.85,2.95) -- ++(-0.13,-0.2) -- ++(0.26,0) -- cycle;
  \fill[ink] (1.4,1.55) -- ++(-0.13,-0.2) -- ++(0.26,0) -- cycle;
  \node[lbl, left] at (0.7,3.0) {A, 11:02};
  \node[lbl, left] at (1.25,1.45) {B, 11:06};
  \draw[track i] plot[smooth] coordinates {(0.85,2.95) (2.6,2.65) (6,3.0) (9,2.65) (10.6,3.0) (12.2,3.9) (13.4,4.3)};
  \draw[track j] plot[smooth] coordinates {(1.4,1.55) (3.2,2.1) (6,2.5) (9,2.15) (12,2.35) (14.6,2.5)};
  \draw[track k] plot[smooth] coordinates {(7.2,0.0) (8.2,1.2) (9.6,1.95) (12,2.15) (14.6,2.25)};
  \node[lbl, blueI, above] at (13.4,4.3) {$i$};
  \node[lbl, orangeJ, right] at (14.6,2.5) {$j$};
  \node[lbl, aqua, right] at (14.6,2.2) {$k$};
  \node[tag] (r1) at (-0.6,0.6) {R1};
  \draw[ink2] (r1) -- (0.5,1.3);
  \node[note, align=left, anchor=north west] at (-1.05,0.3) {same area,\\two takeoffs};
  \node[tag] (r2) at (3.0,0.4) {R2};
  \draw[ink2] (r2) -- (2.0,1.4);
  \node[note, align=left, anchor=north west] at (2.6,0.1) {staggered departures,\\within sight};
  \node[tag] (r3) at (5.6,4.2) {R3};
  \draw[ink2] (r3) -- (6.0,3.15);
  \node[note, align=left, anchor=west] at (6.0,4.2) {nearby places, similar motion};
  \node[tag] (r4) at (10.3,0.5) {R4};
  \draw[ink2] (r4) -- (8.4,1.35);
  \draw[ink2] (r4) -- (11.5,3.45);
  \node[note, align=left, anchor=north west] at (9.9,0.2) {$k$ joins, $i$ changes route};
\end{tikzpicture}
\caption{The four requirements on a schematic map. Pilots $i$ and $j$ take off from two places of
the same area four minutes apart and follow the same corridor (grey) without flying over the same
points; $k$ joins later and $i$ leaves the corridor.}
\label{fig:requirements}
\end{figure}

Three immediate approaches do not meet these requirements (Figure~\ref{fig:approaches}).
\begin{itemize}
  \item \textbf{Geometric similarity of the tracks} (Hausdorff, Fréchet, unconstrained DTW;
    survey in~\cite{su2020}). Cross-country routes are channelled by the terrain and thermals
    are often anchored to the ground: two independent pilots who took off from the same site
    produce similar tracks even hours apart. Unconstrained DTW, moreover, warps the time axis
    without limit and erases exactly the information that (R2) needs.
  \item \textbf{Strict synchrony} (flock~\cite{gudmundsson2006}, convoy~\cite{jeung2008},
    swarm~\cite{li2010}). These patterns require proximity at the same instant; the swarm
    admits non-consecutive instants but no delay between members. Two pilots on the same route
    two minutes apart are never close at the same instant, yet they fly together in the sense
    of (R2)--(R3).
  \item \textbf{Window on the departures.} A rule such as ``departures within $T$ minutes''
    produces groups that depend on the reference pilot: if 1 and 8 depart within the window,
    and so do 8 and 14, then 1 and 14 need not (Section~\ref{sec:transitivity}).
\end{itemize}

\begin{figure}
\centering
\begin{tikzpicture}[x=1cm,y=1cm]
  % (a) geometric similarity
  \begin{scope}[shift={(0,0)}]
    \node[panel] at (0,3.3) {(a) Geometric similarity};
    \draw[muted, line width=0.5pt] plot[smooth] coordinates {(0,2.6) (1.5,2.95) (3,2.7) (4.6,3.0)};
    \draw[muted, line width=0.5pt] plot[smooth] coordinates {(0,0.6) (1.5,0.95) (3,0.7) (4.6,1.0)};
    \node[note, anchor=north] at (3.9,2.75) {valley sides};
    \draw[track i] plot[smooth] coordinates {(0.1,1.6) (1.2,2.0) (2.4,1.7) (3.5,2.05) (4.5,1.85)};
    \draw[track j] plot[smooth] coordinates {(0.1,1.45) (1.2,1.83) (2.4,1.55) (3.5,1.88) (4.5,1.7)};
    \node[lbl, blueI, above] at (0.6,1.85) {$i$, 10:00};
    \node[lbl, orangeJ, below] at (0.75,1.4) {$j$, 14:00};
    \node[note, align=left, anchor=north west] at (0,0.35) {same shape,\\four hours apart};
  \end{scope}
  % (b) strict synchrony
  \begin{scope}[shift={(5.3,0)}]
    \node[panel] at (0,3.3) {(b) Strict synchrony};
    \draw[light, line width=6pt] (0.1,1.5) -- (4.6,1.5);
    \draw[track i] (0.1,1.5) -- (3.4,1.5);
    \draw[track j] (0.1,1.5) -- (1.4,1.5);
    \node[pt i] at (3.4,1.5) {};
    \node[pt j] at (1.4,1.5) {};
    \draw[tol] (3.4,1.5) circle (0.7);
    \draw[<->, ink2] (1.4,2.35) -- node[above, lbl] {2 min $\approx$ 1 km} (3.4,2.35);
    \node[lbl, orangeJ, below=3pt] at (1.4,1.5) {$j$};
    \node[lbl, blueI, below=3pt] at (3.4,1.5) {$i$};
    \node[note, align=left, anchor=north west] at (0,0.35) {together, never close\\at the same instant};
  \end{scope}
  % (c) window on departures
  \begin{scope}[shift={(10.6,0)}]
    \node[panel] at (0,3.3) {(c) Window on departures};
    \draw[->, ink2] (0,1.0) -- (4.9,1.0) node[right, lbl, ink] {$t$};
    \foreach \x/\f/\h in {0.4/1/11:00, 2.0/8/11:08, 3.6/14/11:16} {
      \fill[ink] (\x,1.0) circle (1.6pt);
      \node[lbl, below] at (\x,0.95) {\h};
      \node[lbl, above] at (\x,1.05) {\f};
    }
    \draw[ink2] (0.4,1.6) -- (2.0,1.6) node[midway, above, lbl] {\checkmark};
    \draw[ink2] (2.0,1.95) -- (3.6,1.95) node[midway, above, lbl] {\checkmark};
    \draw[ink2] (0.4,2.45) -- (3.6,2.45) node[midway, above, lbl] {$\times$};
    \node[note, align=left, anchor=north west] at (0,0.35) {window $T = 10$ min: the group\\depends on the reference pilot};
  \end{scope}
\end{tikzpicture}
\caption{Why three immediate rules fail. (a) Two independent pilots on the same valley produce
similar tracks hours apart. (b) Two pilots on the same route two minutes apart never enter each
other's proximity disc. (c) With a window of 10 minutes, 1--8 and 8--14 pass and 1--14 fails.}
\label{fig:approaches}
\end{figure}
